% Transcription — fonds Alexandre Grothendieck, shelfmark 158, batch 1 (pages 1-20).
% Established from the facsimile archives/batches/158/batch-01.pdf (a mirror of
% https://grothendieck.umontpellier.fr/158.pdf), per the skill
% transcribe-grothendieck. No cover sheet: PDF page N = archivists' page N.
% Checked against the pencilled numbers bottom-left, which are legible on
% pages 1, 2, 3, 4, 5, 6, 13, 14, 17, 19 and 20 and agree with the PDF page
% throughout; the batch PDF is twenty pages for twenty leaves, and the folder's
% own PDF would carry the offset that this mirror does not.
%
% Pages 5 and 16 carry nothing but the archivists' pencilled number — absent.
% Every other leaf of the batch is mathematics in his hand and is transcribed.
% Page 7 was scanned upside down and is turned here; page 20 is written on the
% back of a printed letterhead, which is the paper and not the subject.
%
% The batch holds three distinct runs and they are not one argument.
%
% Pages 1 to 4 are four loose sheets of trial: diagrams, formulas and half
% sentences thrown across the leaf, building one object — the category
% pi-tilde\B attached to a functor pi : B --> Ens by composing it with the
% non-empty-parts functor into ordered sets. Page 2's lower half is cancelled
% by two long diagonals and repeats the product computation that page 1 carries.
%
% Pages 6 and 8 to 15 are the run proper, under his own pagination 1 to 9
% standing at the head of pages 6, 8, 9, 10, 11, 12, 13, 14 and 15. Page 7 is
% interpolated between his 1 and his 2, carries no number of his, is struck
% through by two diagonals, and is an abandoned first mouture of the same
% opening — it is kept at its rank because the difference between it and page 6
% is the information. The run starts from a category X such that every
% 0-connected object of X^ is 1-connected, takes F in X^ and the topos induced
% on the image U of F, forms pi_F(I) = pi_0(F^I) on the category Delta-tilde of
% non-empty finite sets, and builds a functor phi from (Delta-tilde/pi)° to Y^
% sending (Delta-tilde_n, alpha) to the connected component alpha of F^{n+1}.
% Extended by continuity, phi is the inverse image of a morphism of topoi
% f : Y^ --> Z^ with Z = (Delta-tilde/pi)°, and pages 9 and 10 show that f is
% W_0-aspheric, then W_1-aspheric under the standing hypothesis, and conclude
% that every 0-connected object of Z^ is 1-connected. Page 10 then conjectures
% the Lemme (?) that would force pi to be the final object, and draws from it
% Corollaire 1 (the cartesian powers F^n of a 0-connected F are 0-connected)
% and Corollaire 2 (X is locally cofiltering). Pages 11 to 15 generalise
% Delta-tilde to an arbitrary A, rebuild pi-tilde\B in that generality, and
% state Lemme 1 on the existence of limits in a fibred category and its
% Corollaire on pi-tilde\B. The run breaks off five lines into page 15 and the
% rest of that leaf is blank.
%
% Pages 17 to 20 are a third run, paginated 0, then 1 and 2 on pages 19 and 20,
% with page 18 interpolated and unnumbered. It is a presentation of the
% fundamental group of a category: page 17 defines the group Delta(E) on the
% brackets [x,y] and states that pi_1(X,e), for X with least object e and
% greatest object a, is the quotient of Delta(Hom(e,a)) by the relations
% [alpha v, alpha u] = [beta v, beta u]. Page 18 is the cancelled scratch sheet
% for exactly that relation, crossed by two diagonals; what survives on it is
% the composition H_{m,n} x H_{n,p} --> H_{m,p} and the question whether
% H_x --> H_y is an isomorphism. Pages 19 and 20 restate the matter for a
% pseudo-cofiltering X with least object e: Proposition 1 on the surjectivity
% of pi_1(X_0,e) --> pi_1(X,e), the well-definedness of [v,u] in G, and the
% kernel of that surjection as the invariant subgroup generated by the elements
% attached to composable triangles — with the N.B. that if e is also the
% greatest object the relation collapses and the map is an isomorphism.
%
% Notation, fixed once for the whole batch rather than page by page. The capital
% P drawn with a crossed bowl is the parts functor and is set \mathfrak{P},
% starred \mathfrak{P}^{*} where the page stars it for the non-empty parts —
% the same glyph and the same reading as folder 104, which writes
% I_n = \mathfrak{P}^*(\Delta_n). It must not be confused with the roman P of
% page 1, the functor B --> X, which is set \mathrm{P}; both stand on page 1
% within four lines of each other. His raised caret after a capital, for the
% category of presheaves, is set X^{\wedge}, Z^{\wedge}, B^{\wedge}, and
% e_{X^{\wedge}} is the final object there. The tilde'd Delta of pages 6 to 15
% is the category of non-empty finite sets Delta-tilde_n = [0,n] inter Z and is
% set \widetilde{\Delta}; the plain Delta of page 17 is a group and has nothing
% to do with it, and the coincidence of letter is his. lim with an arrow drawn
% beneath is set \varprojlim and \varinjlim, per the house convention. E-tilde
% is always the sum over E taken in Z^{\wedge}, and is set \widetilde{E}.
%
% Readings flagged and not settled: the word after « est la cat. » on page 3,
% read « cofibrée » on the sense; the adjective after « un » on page 14's
% Lemme 1, where the ink will not decide between « fibré » and a longer word;
% the subscript of the last pi_1 on page 20, which appears to read X_0 where
% the argument wants X. The diagonal note in the lower-left corner of page 8 is
% given as illegible with the three fragments that do come up under rotation.
% Pages 12, 13 and 15 are the least legible leaves of the batch: the formulas
% carry and the connective prose largely does not.
%
% Pass: Opus 5 (claude-opus-5), 2026-09-20 — first pass, unchecked against the pages by a human.
\documentclass[11pt,a4paper]{article}
% Le préambule commun à toutes les transcriptions du fonds Grothendieck.
%
% Il tient en une idée : **l'incertitude se compose, elle ne se résout pas.**
% Une transcription qui lisse un mot douteux a détruit la seule chose pour
% laquelle on la faisait — un lecteur doit pouvoir savoir, à chaque endroit,
% ce qui était sur le feuillet et ce qui a été ajouté. Les six macros qui
% suivent sont donc l'appareil critique, et non de la décoration.
%
% Les mêmes commandes sont reconnues par `scripts/render.mjs`, qui produit la
% vue de lecture du site. Une macro ajoutée ici doit l'être aussi là-bas,
% faute de quoi le rendu échouera bruyamment — ce qui est voulu : un rendu
% silencieusement faux est pire qu'un rendu absent.

\ProvidesPackage{grothendieck}[2026/08/08 Transcriptions du fonds Grothendieck]

\usepackage{amsmath,amssymb,amsthm}
\usepackage{mathrsfs}
\usepackage{stmaryrd}
% Les diagrammes commutatifs sont partout dans ce fonds : c'est la langue
% même de la géométrie algébrique de Grothendieck, et une transcription qui
% les rendrait en prose ne transcrirait rien.
\usepackage{tikz-cd}
% Français par défaut — c'est la langue des feuillets — mais l'anglais est
% chargé aussi : la traduction et le résumé sont des documents anglais, et la
% typographie française (espaces avant les deux-points) y serait une faute.
% Ces documents basculent par \selectlanguage{english} en tête de corps.
\usepackage[english,french]{babel}
\usepackage{fontspec}
\usepackage{geometry}
\geometry{a4paper,margin=25mm,marginparwidth=28mm}
\usepackage{marginnote}
\usepackage{xcolor}
% ulem, not soul. `soul`'s \st and \ul are built on a character-by-character
% scan that XeLaTeX's font machinery breaks: the document fails with an
% undefined control sequence rather than degrading. ulem does the same two jobs
% and is Unicode-engine safe. `normalem` stops it hijacking \emph, which the
% transcriptions use for Grothendieck's own emphasis.
\usepackage[normalem]{ulem}
\usepackage{hyperref}
\hypersetup{colorlinks=true,linkcolor=black,urlcolor=black,pdfborder={0 0 0}}

% --- Métadonnées du lot -----------------------------------------------------
% Reprises telles quelles par le script de rendu, qui les affiche en tête de
% la vue de lecture. Elles fixent ce que le fichier prétend couvrir : sans
% elles, une transcription flotte sans point d'attache dans un fonds de seize
% mille feuillets.

\newcommand{\folder}[1]{\gdef\grothFolder{#1}}
\newcommand{\batch}[1]{\gdef\grothBatch{#1}}
\newcommand{\leaves}[2]{\gdef\grothFirst{#1}\gdef\grothLast{#2}}
\newcommand{\foldertitle}[1]{\gdef\grothFoldertitle{#1}}
\gdef\grothFolder{?}\gdef\grothBatch{?}\gdef\grothFirst{?}\gdef\grothLast{?}\gdef\grothFoldertitle{}

% --- Le résumé --------------------------------------------------------------
%
% Il ouvre la lecture modernisée, sous le titre « Résumé », et il est composé
% en petit. Ce n'est pas un ornement : c'est le seul passage du document qui
% ne soit pas des mathématiques, et un lecteur doit pouvoir le reconnaître
% avant de l'avoir lu — puis le sauter s'il connaît déjà le sujet, ou s'y
% arrêter s'il ne le connaît pas. Le filet qui le suit marque la reprise du
% corps.

\newenvironment{resume}
  {\small\setlength{\parskip}{0.45em}}
  {\par\medskip\hrule\medskip}

% --- Mots-clés ---------------------------------------------------------------
%
% En anglais, en clôture du résumé de la lecture modernisée : le vocabulaire
% moderne sous lequel chercher ce que les feuillets construisent. Le manifeste
% du site (`npm run manifest`) les extrait pour étiqueter la cote entière —
% cette ligne est la source unique des tags, il n'y a pas de fichier à côté.
\newcommand{\keywords}[1]{\par\smallskip\noindent
  {\footnotesize\textsc{Keywords} --- \textit{#1}}\par}

% --- L'appareil critique ----------------------------------------------------

% Le numéro de feuillet, en marge. C'est l'ancre qui permet de régler un
% désaccord sur une formule en regardant *un* feuillet plutôt qu'une chemise
% de six cents. Le site s'en sert aussi pour tourner les pages du fac-similé
% quand on fait défiler la transcription.
\newcommand{\leaf}[1]{%
  \marginnote{\footnotesize\color{gray!70}\textsf{#1}}%
  \label{leaf:#1}\ignorespaces}

% Illisible : on ne devine pas. Un mot inventé qui se lit comme les autres est
% le pire résultat possible.
\newcommand{\ill}{\textcolor{red!60!black}{[\,\ldots\,]}}

% Lecture douteuse : le mot est donné, et signalé comme incertain.
\newcommand{\uncertain}[1]{\uline{#1}}

% Ajout de l'éditeur — une lettre manquante, un mot sous-entendu.
\newcommand{\add}[1]{\textcolor{blue!50!black}{[#1]}}

% Biffé par Grothendieck lui-même. Ce qu'il a rayé fait partie du document :
% une rature dit souvent où la pensée a changé de direction.
\newcommand{\struck}[1]{\sout{#1}}

% Note du transcripteur, sur l'état matériel du feuillet.
\newcommand{\note}[1]{\par\smallskip{\small\color{gray!60!black}\textit{[#1]}}\par\smallskip}

% Note marginale de Grothendieck, distincte du corps du texte.
\newcommand{\marginal}[1]{\par\smallskip\noindent\hspace{1em}%
  {\small\color{gray!60!black}#1}\par\smallskip}

% --- Titre ------------------------------------------------------------------

\AtBeginDocument{%
  \begin{center}
    {\large\bfseries Fonds Alexandre Grothendieck --- cote n\textsuperscript{o}~\grothFolder}\\[2pt]
    {\itshape\grothFoldertitle}\\[4pt]
    {\small feuillets \grothFirst--\grothLast\ (lot \grothBatch)}\\[2pt]
    {\footnotesize\color{gray!60!black}Transcription établie d'après le fac-similé numérisé
     par l'Université de Montpellier.\\
     Les lectures incertaines sont soulignées, les illisibles notées [\,\ldots\,],
     les ajouts entre crochets.}
  \end{center}
  \vspace{1em}}

\endinput


\folder{158}
\batch{1}
\pages{1}{20}
\dating{[vers 1990-1991]}
\watermark{Édition de démonstration}
\foldertitle{[Autour des Dérivateurs] : notes manuscrites (s.d.)}

\begin{document}

\section*{Feuillets d'essai : la catégorie $\widetilde{\pi}\backslash B$}

\note{les trois titres de sections sont de l'édition : aucun feuillet de ce lot
ne porte de titre de lui}

\page{1}

\note{feuillet d'essai, sans phrase suivie : la prose de liaison est ici
compressée et ce qui ne porte pas est laissé \ill{}}

$X$ topos loc. connexe, $\pi = \pi_{0} \circ \mathrm{P}$ :

\begin{tikzcd}
B \arrow[r, "\mathrm{P}"] \arrow[rr, bend right=35, "\pi"'] & X \arrow[r, "\pi_{0}"] & (\mathrm{Ens})
\end{tikzcd}

\ill{} stable par $\varprojlim$ \ill{} \uncertain{de type $\Phi$} \ill{}
$\Phi$-exact à g.

\[
  \mathrm{Hom}(G, F^{I}) = \mathrm{Hom}\bigl(I, \mathrm{Hom}(G,F)\bigr)
\]

\begin{tikzcd}
Z \arrow[r, "i"] \arrow[dr, "p"'] & C \arrow[r, "\overline{\mathrm{P}}"] \arrow[d, "q"] & X \\
 & B \arrow[u, dashed, "\sigma"'] &
\end{tikzcd}

$\overline{\mathrm{P}}$ \ill{} \uncertain{$\Phi$-exact à g.} ?

$\mathrm{P}(b,E) = \widetilde{E}$, \quad
$\mathfrak{P}\bigl(\pi_{0}(\mathrm{P}(b))\bigr) \simeq$ ens. des \ill{} de
$\mathrm{P}(b)$

\struck{$(b_{i}, E_{i})$}

\[
  \mathrm{P}(b,E) \times \mathrm{P}(b',E')
  = \coprod_{(z,z') \in E \times E'} \widetilde{z} \times \widetilde{z}'
\]

avec $\widetilde{z} \subset \pi(b)$, $\widetilde{z}' \subset \pi(b')$, donc
$\subset \pi(b) \times \pi(b') = \pi(b \times b')$, puis \ill{} $\times\,(E \times E')$.

\page{2}

\begin{tikzcd}
C_{c_{\infty}} \arrow[r, "u'_{i*}"] & C_{b_{i}}
\end{tikzcd}

$c_{\infty} \longrightarrow u'^{\,*}_{i}(c_{i})$ \ill{}

\struck{$\pi(b) \longrightarrow \pi(b')$}

\begin{tikzcd}
E \arrow[r, "u"] \arrow[d] & E' \arrow[d] \\
\mathfrak{P}^{*}(E) \arrow[r] & \mathfrak{P}^{*}(E')
\end{tikzcd}

\[
  u(F) \subset F' \quad \Longleftrightarrow \quad F \subset u^{-1}(F')
\]

\[
  \coprod_{z \in E} z \ \times \coprod_{z' \in E'} z'
  = \coprod_{(z,z') \in E \times E'} z \times z'
\]

\note{produits pris dans $Z^{\wedge}$ ; $z \in E \subset \pi(b)$, la lettre
$\pi$ étant ici écrite sur une lettre biffée}

\note{toute la moitié inférieure du feuillet est annulée par deux longs traits
diagonaux ; ce qui subsiste sous les traits est la même construction du produit,
reprise et abandonnée}

\struck{$z$ au-dessus de $b$, $z'$ au-dessus de $b'$, $b \times b'$ avec ses
deux projections, $z \times z'$, $(b \times b') \times {}$ \ill{} (produits
cartésiens)}

\page{3}

\begin{tikzcd}
B \arrow[r, "\pi"] \arrow[rr, bend right=35, "\widetilde{\pi}"'] & (\mathrm{Ens}) \arrow[r, "\mathfrak{P}^{*}"] & (\mathrm{Ord})
\end{tikzcd}

$C = \widetilde{\pi}\backslash B$, avec un foncteur fidèle vers $B$ :

\[
  \mathrm{Hom}_{C}\bigl((b,E),(b',E')\bigr) \hookrightarrow \mathrm{Hom}_{B}(b,b')
\]

où $E \in \widetilde{\pi}(b) = \mathfrak{P}^{*}(\pi(b))$ et
$E' \in \mathfrak{P}^{*}(\pi(b')) = \widetilde{\pi}(b')$, l'image étant

\[
  \{\, u : b \to b' \ \mid \ \widetilde{\pi}(u)(E) \subset E' \,\}
\]

N.B. $\widetilde{\pi}\backslash B$ est la cat. \uncertain{cofibrée} sur $B$
définie par le foncteur composé \struck{\ill{}}
$\widetilde{\pi} : B \to (\mathrm{Ord}) \to (\mathrm{Cat})$, \ill{}
$i \mapsto c_{i} = (b_{i}, E_{i})$.

$b = \varprojlim b_{i} \longrightarrow b_{i}$

\begin{tikzcd}[column sep=small, row sep=small, nodes={font=\scriptsize}]
c \arrow[r, "\mathrm{cocart}"] \arrow[rr, bend left=30, "v_{i}"] & u_{\infty*}(c) = c_{\infty} \arrow[r, "w_{i}"] & c_{i} = (b_{i},E_{i}) \\
b \arrow[r, "u_{\infty}"] \arrow[rr, bend right=30, "u_{i}"'] & b_{\infty} = \varprojlim b_{i} \arrow[r, "u'_{i}"] & b_{i}
\end{tikzcd}

$w'_{i*}(c_{\infty}) \longrightarrow {}$\uncertain{$\xi_{i}$} dans $C_{b_{i}}$.

\page{4}

\[
  \widetilde{\pi}\backslash B \ \simeq \ B/\widetilde{\pi} \ = \ C,
  \qquad \pi\backslash B \ = \ Z
\]

\note{$Z$ est coétale sur $B$ ; la mention est portée le long de la flèche
verticale}

\[
  C \longrightarrow Z^{\wedge}, \qquad (b,E) \longmapsto \coprod_{z \in E} z
\]

\struck{$\mathfrak{P}(Z_{b})$}

Pour $(b,E) \longrightarrow (b',E')$, avec $E \subset Z_{b} \longrightarrow
Z_{b'} \supset E'$ au-dessus de $u : b \to b'$, on obtient

\[
  u^{E,E'}_{*} : \widetilde{E} \longrightarrow \widetilde{E}',
  \qquad \widetilde{E} = \coprod_{z \in E} z, \qquad z \longmapsto \widetilde{z}'
\]

\begin{tikzcd}[column sep=small, row sep=small, nodes={font=\scriptsize}]
Z \arrow[r, "i"] \arrow[d, "p"'] \arrow[rr, bend left=30, "\mathrm{cocart}"] & C \arrow[r, "\varphi"] \arrow[dl, "q"] & Z^{\wedge} \arrow[d, "p_{!}"] \\
B \arrow[rr, "c_{B}"'] & & B^{\wedge}
\end{tikzcd}

\note{sur le feuillet, $p$ est annoté « cofibré » et $q$ « cofibré \uncertain{et}
fibré »}

\[
  (b,E) \times (b',E') = \bigl(b \times b',\ p_{1}^{*}(E) \cap p_{2}^{*}(E')\bigr)
\]

\section*{Objets $0$-connexes et $1$-connexes ; le morphisme $f : Y^{\wedge} \to Z^{\wedge}$}

\note{à partir d'ici et jusqu'au feuillet 15 court sa propre pagination, de 1 à
9, portée en tête des feuillets 6, 8, 9, 10, 11, 12, 13, 14 et 15 ; le
feuillet 7 s'intercale sans numéro}

\page{6}\note{sa p. 1}

$X \in \mathrm{Ob}\,(\mathrm{Cat})$ tel que tout objet $0$-connexe de
$X^{\wedge}$ soit $1$-connexe. $F \in \mathrm{Ob}\,X^{\wedge}$,
$U = \mathrm{Im}(F \to e_{X^{\wedge}})$. Donc

\[
  F \in \mathrm{Ob}\,X^{\wedge}/U \simeq (X/U)^{\wedge} = Y^{\wedge},
\]

et, plus précisément, $F \to e_{Y^{\wedge}}$ épi.

\note{$Y$ est posé sous l'accolade comme nom de $X/U$}

\struck{\ill{}} Soit $\widetilde{\Delta} = $ cat. des ens. finis de la forme

\[
  \widetilde{\Delta}_{n} = [0,n] \cap \mathbf{Z}
\]

(équivalente à la cat. des ens. finis non vides). Le foncteur
$I \mapsto \pi_{0}(F^{I}) : \widetilde{\Delta}^{\,\circ} \longrightarrow
(\mathrm{Ens})$, soit $\pi_{F} = \pi \in \widetilde{\Delta}^{\wedge}$ :

\[
  \pi(I) = \pi_{0}(F^{I}), \qquad \text{i.e.} \qquad
  \pi(\widetilde{\Delta}_{n}) = \pi_{0}(F^{\,n+1}).
\]

Je vais définir un foncteur

\[
  \varphi : (\widetilde{\Delta}_{/\pi})^{\circ} \longrightarrow Y^{\wedge}
  \simeq (X/U)^{\wedge} \simeq X^{\wedge}/U .
\]

Un objet de $\widetilde{\Delta}_{/\pi}$ est un couple
$(\widetilde{\Delta}_{n}, \alpha)$, où
$\alpha \in \pi(\widetilde{\Delta}_{n}) = \pi_{0}(F^{\,n+1})$ ; on pose

\[
  \varphi(\widetilde{\Delta}_{n}, \alpha) = \alpha \subset F^{\,n+1}
\]

\note{$\alpha$ désigne ici la composante connexe de $F^{\,n+1}$}

\marginal{comme $F$, ses les $F^{I}$ sont des $X^{\wedge}/U \simeq Y^{\wedge}$,
de même de $\alpha$.}

Une flèche de $\widetilde{\Delta}_{/\pi}$, de
$(\widetilde{\Delta}_{n}, \alpha)$ vers $(\widetilde{\Delta}_{m}, \beta)$, est
formée d'une flèche $\rho : \widetilde{\Delta}_{n} \to \widetilde{\Delta}_{m}$
telle que $\rho^{*}_{\pi}(\beta) = \alpha$, où
$\rho^{*}_{\pi}(\beta) \overset{\text{déf}}{=} \pi(\rho)(\beta)$. Or

\note{la flèche est nommée $f$ à sa première mention et $\rho$ partout
ensuite ; la variante est sur le feuillet et n'est pas corrigée ici}

\page{7}

\note{feuillet non numéroté par lui, intercalé entre ses pages 1 et 2, et barré
de deux longs traits diagonaux : c'est une première mouture, abandonnée, du
début repris au feuillet 6}

$X \in \mathrm{Ob}\,(\mathrm{Cat})$ tel que les objets $0$-conn. de
$X^{\wedge}$ soient $1$-connexes. $F \in \mathrm{Ob}\,X^{\wedge}$, d'où
$I \mapsto F^{I}$ sur $\widetilde{\Delta}^{\,\circ}$, les ens. finis non vides,
d'où

\begin{tikzcd}
(\widetilde{\Delta}^{\,\circ})^{\wedge} \arrow[r] & X^{\wedge}
\end{tikzcd}

foncteur qui commute aux $\varinjlim$, avec
$(\widetilde{\Delta}^{\,\circ})^{\wedge} = \mathrm{Hom}(\widetilde{\Delta},
(\mathrm{Ens})) = \widetilde{\Delta}^{\wedge}$.

Soit $U \subset e_{X^{\wedge}}$ l'image de $F \to e_{X^{\wedge}}$ ; alors
$\widetilde{\Delta}^{\,\circ} \longrightarrow X^{\wedge}$ se factorise par
$(X/U)^{\wedge}$, d'où

\[
  \widetilde{\Delta}^{\wedge} \longrightarrow (X/U)^{\wedge} \simeq X^{\wedge}/U
\]

\note{en marge, le morphisme de topos $X^{\wedge}/U \to X^{\wedge}$}

\struck{commute aux $\varinjlim$ \ill{} plus $\varprojlim$}

\begin{tikzcd}
Y^{\wedge} \arrow[r, "f"] & Z^{\wedge}
\end{tikzcd}

\[
  f^{*}(G^{\Psi}) = f^{*}(G)^{\,f^{*}(\Psi)}
\]

\page{8}\note{sa p. 2}

$\rho$ définit $\rho^{*} : F^{\widetilde{\Delta}_{m}} \longrightarrow
F^{\widetilde{\Delta}_{n}}$, et la relation $\rho^{*}_{\pi}(\beta) = \alpha$
signifie que $\rho^{*}|\beta$ se factorise par $\alpha$, i.e. $\rho^{*}$ induit

\[
  \varphi(\rho) : \beta \longrightarrow \alpha,
  \qquad \beta = \varphi(\widetilde{\Delta}_{m}, \beta),
  \qquad \alpha = \varphi(\widetilde{\Delta}_{n}, \alpha).
\]

\note{sur le feuillet, les deux identifications sont portées en doubles barres
verticales sous $\beta$ et sous $\alpha$}

On a ainsi défini une loi fonctorielle $\varphi$. Le foncteur $\varphi$ se
prolonge par continuité en

\[
  \Phi : \bigl((\widetilde{\Delta}_{/\pi})^{\circ}\bigr)^{\wedge}
  \longrightarrow Y^{\wedge},
  \qquad
  \bigl((\widetilde{\Delta}_{/\pi})^{\circ}\bigr)^{\wedge}
  = (\widetilde{\Delta}_{/\pi})^{\wedge} = Z^{\wedge},
\]

foncteur commutant aux $\varinjlim$ quelconques. Ce foncteur est \emph{exact},
d'où qu'il correspond à un morphisme de topos

\[
  f : Y^{\wedge} \longrightarrow Z^{\wedge}
  \qquad \bigl(\text{où } Z \overset{\text{déf}}{=} (\widetilde{\Delta}_{/\pi})^{\circ}\bigr).
\]

\marginal{\ill{} — note portée en diagonale dans le coin inférieur gauche, dont
seuls « c'est ici », « les foncteurs qui interviennent » et « qu'on pourrait »
se laissent lire une fois le feuillet tourné}

\page{9}\note{sa p. 3}

Ce morphisme est $W_{0}$-asphérique, i.e. pour tout $z \in \mathrm{Ob}\,Z$, son
image inverse $f^{*}(z) = \Phi(z)$ ($= \alpha$, si
$z = (\widetilde{\Delta}_{n}, \alpha)$) est $0$-connexe, ce qui implique que
pour tout $G \in \mathrm{Ob}\,Z^{\wedge}$ qui est $0$-connexe,
\struck{\ill{}} $f^{*}(G)$ est $0$-connexe. (Pour le voir, on \add{peut}
utiliser le fait que $f^{*} = \Phi$ commute aux produits fibrés, ou tout au
moins aux intersections de sous-objets.)

Si on fait appel maintenant à l'hyp. que les objets $0$-connexes de $X^{\wedge}$
(donc aussi ceux du topos induit $Y^{\wedge}$) sont $1$-connexes, il s'ensuit
que $f$ est $W_{1}$-asphérique, \emph{plus précisément} : \struck{i.e.} que pour
tout \struck{Pour tout} $G \in \mathrm{Ob}\,Z^{\wedge}$ qui est $0$-connexe,
$f^{*}(G)$ est $1$-connexe.

Mais \uncertain{dans}, c'est un morphisme
$Y^{\wedge}/f^{*}(G) \longrightarrow Z^{\wedge}/G$ \ill{} $0$-connexes,
\ill{} \uncertain{la source} $1$-connexe, \ill{} c'est exact. Donc \ill{} tout
objet $0$-connexe de $Z^{\wedge}$ est $1$-connexe.

\page{10}\note{sa p. 4}

Mais je conjecture que l'on devrait avoir ceci :

\textbf{Lemme (?)} Soient $\pi \in \mathrm{Ob}\,\widetilde{\Delta}^{\wedge}$,
$Z = (\widetilde{\Delta}_{/\pi})^{\circ}$ \ill{}, et supposons que

\begin{itemize}
\item[a)] tout objet $0$-connexe de $Z^{\wedge}$ soit $1$-connexe,
\item[b)] $\pi(\widetilde{\Delta}_{0})$ est ponctuel ;
\end{itemize}

alors $\pi \simeq e_{\widetilde{\Delta}^{\wedge}}$, i.e.
$\pi_{n} = \pi(\widetilde{\Delta}_{n})$ est un pt pour tout $n \in \mathbf{N}$.

Vu ce qui précède, cela implique que

\textbf{Corollaire 1} Soit $X$ un objet de $(\mathrm{Cat})$ tel que tout objet
de $X^{\wedge}$ qui est $0$-connexe soit $1$-connexe. Alors pour tout objet
$0$-connexe $F$ de $X^{\wedge}$, les produits cartésiens $F^{n}$ ($n \geq 1$)
soient $0$-connexes, en particulier $F \times F$ est $0$-connexe.

Vu \struck{le fait} qu'on sait déjà que \struck{l'hyp. implique} $X^{\wedge}$
est loc. indécomposable (i.e. $X$ est localement \uncertain{PS 1}), on
\uncertain{voit} que $X^{\wedge}$ \ill{} $0$-connexe, i.e. $X$ loc.
$0$-connexe, donc

\textbf{Corollaire 2} L'hyp. du cor. 1 implique que $X$ est loc. cofiltrant
(et inversement, bien \ill{}), i.e. les comp. connexes \ill{} sont cofiltrantes.

\page{11}\note{sa p. 5}

$A \in \mathrm{Ob}\,(\mathrm{Cat})$ \note{il ajoute : « je pense en ce $A = \widetilde{\Delta}$ »},
$\pi \in \mathrm{Ob}\,A^{\wedge}$,

\[
  Z = (A_{/\pi})^{\circ} \simeq \pi\backslash A^{\circ} = \pi\backslash B,
  \qquad \pi : A^{\circ} \longrightarrow (\mathrm{Ens}), \quad B = A^{\circ}.
\]

On s'intéresse à $Z^{\wedge}$, et à la \uncertain{sous-}catégorie pleine de
$Z^{\wedge}$ engendrée par les $\varprojlim$ finies, engendrée par \ill{} $Z$.
J'ai vu que les candidats sont ceux-ci : considérons
$Z_{b} \simeq \pi(b)$ pour tout $b \in \mathrm{Ob}\,B$, et pour toute partie non
vide $E$ de $\pi(b)$, $E \in \mathfrak{P}^{*}(\pi(b))$, considérons l'élément

\[
  \widetilde{E} = \coprod_{a \in E} a
  \qquad \text{(somme prise dans } Z^{\wedge}\text{)}.
\]

Si $u : b \longrightarrow b'$ dans $B$, \struck{on a}

\[
  u_{*} = \pi(u) : Z_{b} \longrightarrow Z_{b'},
  \qquad Z_{b} = \pi(b), \quad Z_{b'} = \pi(b'),
\]

et si $E \subset Z_{b}$, $E' \subset Z_{b'}$ sont tels que
$u_{*}(E) \subset E'$, alors $u$ induit

\[
  u^{E,E'}_{*} : \widetilde{E} \longrightarrow \widetilde{E}' .
\]

Ainsi, le foncteur composé
$\widetilde{\pi} : b \longmapsto \mathfrak{P}^{*}(\pi(b))$,

\begin{tikzcd}
B \arrow[r, "\pi"] & (\mathrm{Ens}) \arrow[r, "\mathfrak{P}^{*}"] & (\mathrm{Ord})
\end{tikzcd}

définit une catégorie $\widetilde{\pi}\backslash B$, \struck{dont} \ill{}

\page{12}\note{sa p. 6}

$\widetilde{\pi}\backslash B$, dont les objets \ill{} (i.e. $\widetilde{\pi}$
est le foncteur composé de $\widetilde{\pi} : B \to (\mathrm{Ord})$ \ill{} et
$(\mathrm{Ord}) \to (\mathrm{Ens})$), mais \ill{} des flèches : une flèche de
$(b,E)$ dans $(b',E')$ (où $E \in \widetilde{\pi}(b)$,
$E' \in \widetilde{\pi}(b')$) est une flèche $u : b \to b'$ telle que

\[
  \widetilde{\pi}(u)(E) \subset E' \qquad \text{(et non } \widetilde{\pi}(u)(E) = E').
\]

$(*)$ foncteur canonique

\[
  \widetilde{\pi}\backslash B \longrightarrow (\pi\backslash B)^{\wedge} = Z^{\wedge},
  \qquad
  (b,E) \longmapsto \widetilde{E} = \coprod_{a \in E} a \ \subset
  \coprod_{a \in Z_{b}} a,
\]

où $E \in \mathfrak{P}^{*}(\pi(b))$ et $Z_{b} = \pi(b)$.

Je voudrais des conditions sur $A$, ou encore sur $B = A^{\circ}$ et
$\pi \in \mathrm{Ob}\,A^{\wedge} = \mathrm{Ob}\,B^{\wedge}$, qui assurent que
\struck{\ill{}} le foncteur $(*)$ soit \emph{pleinement fidèle}, \ill{} et
\ill{} pour les $\varprojlim$ finies. Par ailleurs, il est clair que \ill{}
l'image essentielle contient $Z$. J'imagine qu'il \ill{} déraisonnable
\supplied{de s'}attendre \ill{}.

\page{13}\note{sa p. 7}

\ill{} pleinement fidèle. Je vais prendre \uncertain{exigeons} que
\struck{$\mathrm{Ob}$} $\widetilde{\pi}\backslash B$ soit stable pour les
$\varprojlim$ finies, et que le foncteur \ill{} soit \ill{}.

Prenons pour $\pi$ le foncteur constant de valeur $e$ ; on trouve alors que

\[
  \widetilde{\pi}\backslash B \ \simeq \ \pi\backslash B \ \simeq \ B,
\]

\ill{} et les \ill{} de $B$ sont stables par $\varprojlim$ \ill{} non vides
\ill{}, $A$ stables par $\varprojlim$ finies, i.e. \ill{} pour $B$, cela revient
aux stabilités

\begin{itemize}
\item[a)] par produit binaire, et
\item[b)] par produit fibré binaire (ou, \uncertain{au} choix, \ill{} par
noyaux des couples de flèches),
\end{itemize}

qui alors $\widetilde{\pi}\backslash B$ est stable par les \ill{} opérations.

En fait, il \uncertain{faudra} faire une hyp. plus \ill{} sur $\pi$, comme que
$C = \widetilde{\pi}\backslash B$ soit non seulement cofibré sur $B$, mais
\ill{} soit \emph{fibré}. Cette hypothèse \ill{} plusieurs \ill{} de $B$ \ill{}

\page{14}\note{sa p. 8}

\ill{} éclaircissement. Alors, pour $u : b \longrightarrow b'$ donné,
l'\struck{application} application

\[
  u_{*} = \widetilde{\pi}(u) : \widetilde{\pi}(b) \longrightarrow
  \widetilde{\pi}(b'),
  \qquad
  \mathfrak{P}^{*}(\pi(b)) \longrightarrow \mathfrak{P}^{*}(\pi(b')),
  \qquad E \longmapsto \pi(u)(E),
\]

\note{le signe $=$ de la première formule est écrit sur un symbole biffé}

admet un adjoint à droite $u^{*}$, à savoir
$E' \longmapsto \pi(u)^{-1}(E')$, qui transforme \ill{} non vides en non vides
(\uncertain{second} \ill{} si $\pi(b) \longrightarrow \pi(b')$ surjective).

\textbf{Lemme 1} Soit $C \longrightarrow B$ un \uncertain{fibré} \ill{} dans
$(\mathrm{Cat})$, \struck{et} soit $I$ dans $(\mathrm{Cat})$ telle que les
$\varprojlim$ du type $I$ existent dans $B$, et dans les fibres $C_{b}$. Alors
elles existent dans $C$, et $C \longrightarrow B$ commute aux $\varprojlim$ de
type $I$.

\textbf{Corollaire} Soit $\pi : B \longrightarrow (\mathrm{Ens})$ un foncteur,
\struck{\ill{}} soit

\begin{tikzcd}
B \arrow[r, "\pi"] & (\mathrm{Ens}) \arrow[r, "\mathfrak{P}^{*}"] & (\mathrm{Ord}) \arrow[r, "\mathrm{can.}"] & (\mathrm{Cat})
\end{tikzcd}

le composé $\widetilde{\pi}$. Alors $C = \widetilde{\pi}\backslash B$ est \ill{}

\page{15}\note{sa p. 9 ; le texte s'arrête au bout de cinq lignes et le reste du feuillet est blanc}

\ill{} catégorie cofibrée sur $B$, à fibres les
$\widetilde{\pi}(b) = {}$\struck{$(\mathrm{Cat})$} $\mathfrak{P}^{*}(\pi(b))$,
définie par $\widetilde{\pi}$, et \ill{} comme les fibres sur $B$,
\ill{} d'\uncertain{abord} par $\widetilde{\pi}$, \ill{} $B$ \ill{}
(resp. induites) \ill{} $I$ \struck{telle que} \ill{} type de $\varprojlim$
\struck{\ill{}} $=$ est stable.

\section*{Une présentation de $\pi_{1}$ d'une catégorie}

\page{17}\note{sa p. 0}

\[
  \Delta(E) = \text{groupe engendré par les } [x,y], \quad (x,y) \in E^{2},
\]

avec les relations

\[
  [x,x] = 1, \qquad [x,y][y,x] = 1, \qquad [x,y][y,z] = [x,z].
\]

\note{ce $\Delta$ est un groupe et n'a rien à voir avec le
$\widetilde{\Delta}$ des feuillets 6 à 15 ; la coïncidence de lettre est de lui}

N.B. Je \uncertain{suppose} pas que $\Delta(E)$ est isomorphe au sous-groupe des
groupes libres $L(E)$ engendré par $E$, engendré par les éléments de la forme
$[x,y] = x^{-1}y$.

Si $\mathrm{card}\,E = n$ est fini, $\Delta(E)$ est le groupe libre à \ill{}
générateurs ; si $E = [1,n]$, il \struck{\ill{}} est engendré par les
générateurs libres \struck{$E_{i}$}

\[
  g_{i} = [i, i+1] \qquad \text{pour } i \in [1, n-1].
\]

\note{il accompagne la formule d'une rangée de points numérotés
$1$, $2$, $3$, \ill{}, $n-1$, $n$}

On a alors, pour $i < j$,

\[
  \begin{cases}
    [i,j] = g_{i}\,g_{i+1} \cdots g_{j-1} \\
    [j,i] = [i,j]^{-1} = g_{j-1}^{-1} \cdots g_{i+1}^{-1} g_{i}^{-1}
  \end{cases}
\]

\uncertain{Résulte} \struck{\ill{}} que si $E$ est dénombrable alors
$\Delta(E)$ est un groupe libre avec un ensemble dénombrable de générateurs
libres.

Soit $X$ une catégorie ayant un plus petit objet $e$ et un plus grand objet
$a$. Alors $\pi_{1}(X,e)$ est isomorphe \struck{\ill{}} au quotient du groupe
$\Delta(\mathrm{Hom}(e,a))$ par les relations suivantes : pour tout système de
doubles flèches

\begin{tikzcd}
e \arrow[r, bend left=25, "u"] \arrow[r, bend right=25, "v"'] & x \arrow[r, bend left=25, "\alpha"] \arrow[r, bend right=25, "\beta"'] & a
\end{tikzcd}

on a la relation

\[
  [\alpha v, \alpha u] = [\beta v, \beta u],
  \qquad \text{i.e.} \qquad [\alpha v, \alpha u]\,[\beta v, \beta u]^{-1} = 1 .
\]

\page{18}

\note{feuillet non numéroté par lui, barré de deux longs traits diagonaux : ce
sont les essais pour la relation du feuillet 17, et seuls les fragments qui
portent sont donnés ici}

\struck{$x = (a_{i})_{i \in I}$} \qquad \struck{$X/a$} \qquad
\struck{$X//a_{i}$} \qquad \struck{$X//a$}

\begin{tikzcd}
z \arrow[r] & x \arrow[r, bend left=25] \arrow[r, bend right=25] & y
\end{tikzcd}

\begin{tikzcd}
& a_{j} \\
a \arrow[ur] \arrow[dr] & \\
& a_{i}
\end{tikzcd}

\begin{tikzcd}
x_{1} \arrow[r, bend left=25] \arrow[r, bend right=25] & x_{0}
\end{tikzcd}

\note{une seconde paire parallèle, de même forme mais dont les deux buts ne se
lisent pas, est tracée au-dessus de celle-ci, ses flèches nommées $u'$, $v'$
puis $u$, $v$}

\[
  H_{m,n} \times H_{n,p} \longrightarrow H_{m,p}
\]

\[
  fu = gv, \qquad gv' = gw, \qquad fuv_{1} = gwv_{1} = gv'v'_{1} = gw\,v'_{1}
\]

\[
  H_{x} = x/R_{x}
\]

\begin{tikzcd}
x \arrow[r] \arrow[d] & y \arrow[d] \\
H_{x} \arrow[r] \arrow[d, "\mathrm{iso}?"] & H_{y} \\
\lbrack x \rbrack &
\end{tikzcd}

\page{19}\note{sa p. 1}

\textbf{Proposition 1} Soit $X \in \mathrm{Ob}\,(\mathrm{Cat})$ une catégorie
\struck{\ill{}} pseudo-cofiltrante ($\Longleftrightarrow X^{\wedge}$ est
localement irréductible), ayant un plus petit objet $e$ (donc $X$ est
$0$-connexe). Soit $X_{0}$ la sous-catégorie pleine réduite à $e$,

\[
  M = \mathrm{End}_{X}(e) = \mathrm{End}_{X_{0}}(e)
\]

le monoïde des endomorphismes de $e$, et
$G = \pi_{1}(M) = \pi_{1}(X_{0}, e)$ le groupe \uncertain{enveloppant} de $M$.
Alors

\begin{itemize}
\item[a)] L'homomorphisme canonique
$\pi_{1}(X_{0}, e) \longrightarrow \pi_{1}(X,e)$ est surjectif, \struck{et pour}
\item[b)] Pour toute double flèche de source $e$
\end{itemize}

\begin{tikzcd}
e \arrow[r, bend left=25, "u"] \arrow[r, bend right=25, "v"'] & x
\end{tikzcd}

dans $X$, il existe $u', v' \in M$ tels que

\[
  uu' = vv',
\]

et l'élément $[v][u]$ de $\pi_{1}(X,e)$ (image canonique des éléments
$e \xrightarrow{\ u\ } x \xleftarrow{\ v\ } e$ dans $\pi_{1}(X,e)$) est égale à
l'image de l'élément

\[
  [v,u] \overset{\text{déf}}{=} v'u'^{-1} \in G = \pi_{1}(X_{0}, e).
\]

De plus, \struck{$e$ est aussi objet maximal de $X$,} cet élément $[v,u]$ ne
dépend pas des choix des couples $(u',v')$ \struck{\ill{}} d'éléments de $M$
satisfaisant $uu' = vv'$. Plus précisément, si $\alpha : x \longrightarrow e$,
alors $u_{1} = \alpha u$ et $v_{1} = \alpha v$ sont des éléments de $M$, et on a
\uncertain{$u_{1}u' = v_{1}v'$}, d'où

\[
  [v,u] = v_{1}^{-1} u_{1} .
\]

\begin{itemize}
\item[c)] Pour toute double-flèche $e \rightrightarrows x$ dans $X$, \ill{}
\end{itemize}

\page{20}\note{sa p. 2}

\ill{} diviseurs $u', v'$ dans $M$, tels que $uu' = vv'$, d'où
$[u,v] \in G$. \add{Alors} le noyau de l'homomorphisme surjectif

\[
  G = \pi_{1}(X_{0}, e) \longrightarrow \pi_{1}(X, e)
\]

est le \uncertain{sous-groupe invariant} de $G$ engendré par les éléments
suivants, correspondants aux diagrammes (\uncertain{par vice-versa}
nécessairement commutatifs)

\begin{tikzcd}
x \arrow[r, "f"] & y \arrow[r, "g"] & z \\
& e \arrow[ul, "u"] \arrow[u, "v"] \arrow[ur, "w"'] &
\end{tikzcd}

\[
  [w, gv]\,[v, fu]\,[w, gfu]^{-1} .
\]

\note{i.e. pour $u', v', w' \in M$ \ill{}}

N.B. S'il existe $\alpha : z \longrightarrow e$, on voit aussitôt que cette
relation est triviale. Donc si $e$ est aussi plus grand objet de $X$, alors

\[
  \pi_{1}(X_{0}, e) \ \xrightarrow{\ \sim\ }\ G
  \ \xrightarrow{\ \sim\ }\ \pi_{1}(X,e) \qquad \text{iso !}
\]

\note{l'indice du premier $\pi_{1}$ se lit $X_{0}$ là où l'argument demanderait
$X$ ; la lecture est douteuse et n'est pas corrigée ici}

Soient $M$ un \add{pseudo-}\ill{} monoïde, $R = R_{\varphi}$ la relation
d'équivalence dans $M$ :

\[
  u \sim_{R} v \quad \Longleftrightarrow \quad \exists\, p, q \in M
\]

\ill{}

(cette relation d'équivalence \ill{} à droite, donc \ill{}
$\overset{\text{déf}}{=} M/R$), \ill{} $M$ \ill{}

\[
  Y = B_{M}, \qquad \varphi \in M, \qquad \varphi f u = \varphi g v ,
\]

\note{il écrit « coïncidences » au-dessus de $M$, et « stable par
translation » sous la formule}

\[
  X = Y/F .
\]

\end{document}
